\documentclass[11pt]{article}
\usepackage[utf8]{inputenc}
\usepackage[T1]{fontenc}
\usepackage[margin=1in]{geometry}
\usepackage{hyperref}
\usepackage{enumitem}
\hypersetup{colorlinks=true,linkcolor=black,urlcolor=blue,citecolor=blue}
\title{Digital Coloring and Tool Routines}
\author{Piter Garcia}
\date{March 20, 2026}
\begin{document}
\maketitle
\section*{Packet Use}
This printable lesson packet is generated from the paired Markdown guide. The Markdown version remains the clickable, neurodivergent-friendly working copy; this PDF carries the same lesson substance in a formal handout format. Evidence claims in this packet are based on transcript-derived lesson notes and the cleaned transcript files in this workspace.

\section*{Quick Scan}
\begin{itemize}[leftmargin=*]
\item \textbf{Grade:} Kindergarten
\item \textbf{Grade Hub:} Notes [../kindergarten-Notes.md]
\item \textbf{Schedule Reference:} Spring2026\_TeachingSchedule\_Derek.md [../../school/Spring2026\_TeachingSchedule\_Derek.md]
\item \textbf{Workbook Sheet:} \texttt{K} in \texttt{Teaching\_Placement-IGNITE Curriculum\_ GCSD25-26.xlsx}
\item \textbf{Lesson Focus:} Device-opening routines, basic digital tool use, color selection, and early classroom computing habits.
\item \textbf{CS Alignment Strength:} Foundational CS / digital fluency
\item \textbf{LaTeX Source:} digital-coloring-and-tool-routines.tex [./digital-coloring-and-tool-routines.tex]
\item \textbf{Bib File:} digital-coloring-and-tool-routines.bib [./digital-coloring-and-tool-routines.bib]
\item \textbf{Artifacts Folder:} artifacts/ [./artifacts/]
\end{itemize}

\section*{Lesson Identity}
\begin{itemize}[leftmargin=*]
\item \textbf{Likely Class Blocks:} Day 1 -- 12:25-1:15 -- Justinger, Day 2 -- 12:25-1:15 -- Kesys, Day 4 -- 12:25-1:15 -- Pum
\item \textbf{Main Lesson Family:} Digital Coloring and Tool Routines
\item \textbf{Primary Evidence Base:} Transcript-derived notes plus source transcripts from this teaching-placement workspace.
\item \textbf{Primary External Source Type:} Official curriculum/provider materials.
\end{itemize}

\section*{What We Are Teaching}
\begin{itemize}[leftmargin=*]
\item opening and closing a classroom app safely
\item choosing simple digital tools such as color, fill, or erase
\item following a repeated click sequence in order
\item sharing attention between teacher model, device, and peer space
\end{itemize}

\section*{CS Standards Alignment}
\begin{itemize}[leftmargin=*]
\item \textbf{1A-CS-02}: Use appropriate terminology to describe common physical and digital parts of a computing system. Why it fits here: Students learn names and functions for screen, click/tap, color tool, and classroom device parts.
\item \textbf{1A-AP-08}: Model daily processes by creating and following algorithms. Why it fits here: Students follow a short repeatable routine: open tool, choose color, complete picture, save or show.
\item \textbf{1A-IC-17}: Work respectfully and responsibly with others online or when using shared devices. Why it fits here: Students practice turn-taking, shared attention, and safe use of common classroom technology.
\end{itemize}

\section*{NYS / Local Curriculum Alignment}
\begin{itemize}[leftmargin=*]
\item Aligned to the local IGNITE kindergarten technology sequence on the \texttt{K} workbook sheet, especially device routines and beginning digital creation.
\item Supports New York-facing digital fluency goals through tool navigation, classroom technology expectations, and simple artifact creation.
\item Functions as a pre-CS lesson that prepares students for later algorithmic work by stabilizing device routines and attention patterns.
\end{itemize}

\section*{Lesson Content and Concepts}
\begin{itemize}[leftmargin=*]
\item device routines
\item tool selection
\item sequence
\item digital artifact
\item shared technology expectations
\end{itemize}

\section*{Evidence / Proof of Instruction}
\begin{itemize}[leftmargin=*]
\item \textbf{Transcript-derived note:} 260318 Kindergarten Justinger Coloring Robotics [../../transcript-derived-lessons/260318-kindergarten-justinger-coloring-robotics.md] The lesson record ties digital coloring to first-week tool routines and introductory robotics vocabulary rather than free play.
\item \textbf{Source transcript:} 260318 Kindergarten Justinger Coloring Robotics [../../transcripts/260318-kindergarten-justinger-coloring-robotics.md] The transcript shows teacher-directed setup, repeated modeling, and support for students who needed help staying with the device task.
\item \textbf{Observed teacher moves:} Instruction is evidenced by modeled clicks, repeated redirects, naming of tools, and circulation for one-to-one support.
\item \textbf{Likely student proof:} Completed digital coloring pages, correct tool use, and students independently re-entering the routine are evidence that the lesson content was taught.
\end{itemize}

\section*{Assessment Evidence}
\begin{itemize}[leftmargin=*]
\item Student can open the target activity and stay in the correct workspace.
\item Student can choose at least one tool intentionally and explain what it does.
\item Student follows a 2-3 step routine with fewer prompts by the end of the lesson.
\end{itemize}

\section*{UDL and Inclusion Supports}
\subsection*{Multiple Representation}
\begin{itemize}[leftmargin=*]
\item Model the same step on the projector, on a student device, and with gesture language.
\item Use visual icon language for tap, color, erase, and done.
\item Keep vocabulary concrete and paired with physical pointing.
\end{itemize}

\subsection*{Multiple Action / Expression}
\begin{itemize}[leftmargin=*]
\item Allow mouse, touchscreen, or partner-supported input.
\item Chunk the task into micro-goals: open, choose, color, show.
\item Accept pointing and verbal naming as partial demonstration of understanding.
\end{itemize}

\subsection*{Multiple Engagement}
\begin{itemize}[leftmargin=*]
\item Use short cycles with fast success feedback.
\item Build in pair praise and quick teacher celebration for correct routines.
\item Keep the task visually appealing and low-text.
\end{itemize}

\section*{How We Are Already Making It Inclusive}
\begin{itemize}[leftmargin=*]
\item Transcript evidence suggests repeated whole-group modeling before release.
\item Teacher circulation and immediate redirection reduced overload for students who lost the routine.
\item The task used low-language, high-visual entry points appropriate for early learners.
\end{itemize}

\section*{How to Improve Inclusion Next Time}
\begin{itemize}[leftmargin=*]
\item Add picture step cards at tables so students can self-recover without waiting.
\item Prepare a sensory-light version with fewer on-screen choices for students who become overstimulated.
\item Create bilingual tool cards and a “help me” nonverbal signal.
\item Capture one finished exemplar in \texttt{artifacts/} as a visual success target.
\end{itemize}

\section*{Materials and Setup}
\begin{itemize}[leftmargin=*]
\item Charged student devices
\item Digital coloring or drawing app
\item Projected teacher model
\item Optional stylus or adaptive pointing tool
\end{itemize}

\section*{Teacher Moves / Script / Routines}
\begin{itemize}[leftmargin=*]
\item Open with a one-sentence goal and show the exact first click.
\item Name each tool only when students need it, not all at once.
\item Release students in short intervals and pause to re-model when multiple students miss the same step.
\item Close by having students show the tool they used or name the routine they remember.
\end{itemize}

\section*{Common Errors and Debugging}
\begin{itemize}[leftmargin=*]
\item Students may tap the wrong icon and leave the activity.
\item Students may color randomly without understanding tool purpose.
\item Some students may freeze when asked to manage both the device and oral directions.
\end{itemize}

\section*{Extensions and Simplifications}
\begin{itemize}[leftmargin=*]
\item Extension: add a label, shape, or robot-related vocabulary word to the image.
\item Simplification: reduce the activity to one color tool and one completion target.
\item Extension: pair the coloring task with a quick talk about what a robot can and cannot do.
\end{itemize}

\section*{Related Local Materials}
\begin{itemize}[leftmargin=*]
\item No direct local lesson packet linked yet.
\end{itemize}

\section*{Artifacts Checklist}
\begin{itemize}[leftmargin=*]
\item Compiled lesson PDF in \texttt{artifacts/}
\item At least one screenshot, student example, or setup photo
\item One short assessment or observation record
\item Updated standards/source bibliography once the \texttt{.bib} file is revised
\item Any teacher reflection or revision notes after reteaching
\end{itemize}

\section*{Selected Official Sources}
Selected official sources that informed this packet are listed below.
ocite{*}

\bibliographystyle{plain}
\bibliography{digital-coloring-and-tool-routines}
\end{document}
